\chapter{Giới thiệu đề tài}
\label{ch:gioithieu}

Các dịch vụ dựa trên vị trí (LBS) đã trở thành một phần không thể thiếu trong đời
sống số: điều hướng, tìm kiếm địa điểm, gọi xe, giao đồ ăn, quảng cáo cá nhân hóa
và các ứng dụng y tế. Sự tiện lợi này đến từ việc người dùng liên tục báo cáo vị
trí của mình cho nhà cung cấp dịch vụ. Tuy nhiên, dữ liệu quỹ đạo --- chuỗi vị trí
theo thời gian --- có thể tiết lộ những thông tin nhạy cảm bậc nhất về một cá
nhân: nơi ở, nơi làm việc, thói quen sinh hoạt, tình trạng sức khỏe, quan hệ xã
hội và niềm tin tôn giáo. Khi dữ liệu này bị lạm dụng, hậu quả trải từ theo dõi,
phân biệt đối xử cho tới đe dọa an toàn cá nhân.

Trong giai đoạn Thực tập 2, chúng tôi đã xây dựng một pipeline bảo vệ quỹ đạo
dựa trên Geo-Indistinguishability (Geo-I)~\cite{andres2013geo}: thêm nhiễu Laplace
phẳng vào từng điểm, giới hạn bán kính nhiễu trong ngưỡng QoS, loại bỏ các điểm
rơi vào tòa nhà/mặt nước, và làm khớp điểm giả vào lưới đường. Khảo sát văn liệu
sâu hơn (Chương~\ref{ch:lienquan}) cho thấy pipeline này tồn tại hai lỗi hình thức
nghiêm trọng: \emph{(i)} việc chặn bán kính nhiễu quanh vị trí thật làm hai vị trí
đủ xa nhau có support đầu ra rời nhau, khiến đảm bảo Geo-I thuần không còn đúng;
\emph{(ii)} việc lấy mẫu lại khi điểm giả rơi vào vùng cấm điều kiện hóa phân phối
đầu ra trên vị trí thật, đồng thời chính loại ``lọc theo tính hợp lý'' này là bề
mặt mà các tấn công tái tạo hiện đại như RAoPT~\cite{buchholz2022raopt} khai thác.

Luận văn này kế thừa mục tiêu của Thực tập 2 nhưng thay thế cơ chế lõi bằng hai cơ
chế mới hoạt động \emph{ngay trên lưới đường}:

\begin{itemize}
  \item \textbf{REM} (Road-network Exponential Mechanism): cơ chế lũy thừa trên
  tập đỉnh của đồ thị đường OSM, thỏa $\eps$-Geo-I hình thức, và mọi vị trí công
  bố đều nằm trên đường \emph{theo cấu trúc} --- loại bỏ hoàn toàn bề mặt tấn công
  map-matching.
  \item \textbf{T-REM}: REM bổ sung trọng số khả đạt (reachability) chỉ phụ thuộc
  vào \emph{điểm đã công bố trước đó} --- thông tin công khai --- nên đảm bảo
  riêng tư per-point không đổi, trong khi chuỗi công bố trở nên nhất quán về tốc
  độ, đóng tấn công liên kết vận tốc (velocity-linkage).
\end{itemize}

Đóng góp của luận văn tính đến bản cập nhật này:
\begin{enumerate}
  \item Phân tích hình thức các lỗi của pipeline Thực tập 2
  (Mục~\ref{sec:phantichbaseline});
  \item Hai cơ chế REM và T-REM kèm chứng minh đảm bảo riêng tư
  (Mục~\ref{sec:rem}, \ref{sec:trem});
  \item Bộ khung đánh giá theo chuẩn văn liệu: tấn công suy luận Bayes tối ưu và
  tấn công theo dõi HMM khai thác tương quan thời gian, cùng các độ đo tiện ích
  và tính thực tế (Chương~\ref{ch:thucnghiem});
  \item Thực nghiệm trên dữ liệu GPS thật (GeoLife, Bắc Kinh) trên đồ thị đường
  OSM thật, so sánh với Laplace phẳng và baseline Thực tập 2;
  \item Hệ thống mô phỏng trực quan ba góc nhìn (người dùng / máy chủ LBS / kẻ
  tấn công) minh họa các kịch bản sử dụng đời thực.
\end{enumerate}

Phần còn lại của luận văn được tổ chức như sau. Chương~\ref{ch:cosolythuyet}
trình bày cơ sở lý thuyết. Chương~\ref{ch:lienquan} khảo sát các công trình liên
quan và xác định khoảng trống nghiên cứu. Chương~\ref{ch:phuongphap} mô tả và
chứng minh các cơ chế đề xuất. Chương~\ref{ch:thucnghiem} trình bày thực nghiệm.
Chương~\ref{ch:tongket} tổng kết và nêu hướng phát triển.
